\section{Ownership: Data, Models, and Weights as User-Owned Assets}
\label{sec:ownership}

The execution proofs make a stronger claim possible than ``the network computed
correctly.'' They make \emph{sovereignty} auditable: a user can own their data, own
the inference performed on it, and own the model weights they improve, and can choose
\emph{per request} how much trust to place in whoever runs the compute---without the
choice ever weakening the guarantee that the work was genuine. Trust becomes a UX
dial, not a security compromise.

\subsection{The trust spectrum}

For any request, a user selects one of three execution modes, ordered by decreasing
sovereignty and increasing convenience. Every mode emits the same execution proofs of
Sections~\ref{sec:four-proofs}--\ref{sec:freivalds}, so every mode is auditable; the
modes differ only in \emph{who holds the data in the clear} and \emph{who the user
must trust for liveness}.

\begin{description}[leftmargin=0pt,style=nextline]
\item[(1) User-operated node --- maximal sovereignty.] The user runs the inference on
their own hardware. The data never leaves the device; the model and weights are local;
the proof is emitted for the user's own records or for third parties who later wish to
verify a result the user publishes. Trust in any operator is zero because there is no
operator. This is the floor of the spectrum and the reference against which the others
are measured.
\item[(2) DAO trustless nodes --- verify, don't trust.] The user sends work to
permissionless DAO operators, but under confidential compute (T1/T2) so the operator
is \emph{plaintext-blind}: it processes ciphertext inside an enclave and never sees the
data in the clear, while the execution proof (\PoI{}/\PoT{}) plus the slashing bond
makes a wrong or fabricated result detectable and punishable. The user trusts no
particular operator's honesty---only the proof system and, for the TEE tiers, the
hardware root. This is the ``verify, don't trust'' mode: sovereignty over correctness
and confidentiality without running the hardware oneself.
\item[(3) Selected trusted operators --- convenience.] The user picks specific
operators by reputation (via the identity-and-reputation layer), accepting a
centralized-feeling, faster, cheaper path. Even here the proofs remain emittable, so
the user retains an audit trail: a trusted operator that misbehaves can be caught
after the fact by anyone who checks the proof, and loses reputation and future
routing. Trust is placed deliberately and remains accountable.
\end{description}

The three modes are points on a continuum, not separate products, because they share
the binding and the proof primitive. A user can move a workload from mode (3) to
mode (2) to mode (1) as their trust requirements rise, with no change to the artifacts
the workload produces.

\subsection{Data is never plaintext at rest}

In modes (2) and (3) the data is encrypted under keys held in the user's key
management, committed on-chain only by hash, and---in the confidential-compute
tiers---processed as ciphertext inside the enclave. The execution proof binds to
$\mathsf{inputHash}$ (Section~\ref{sec:binding}), so the proof certifies that the
\emph{committed} input was processed, while revealing nothing about the input itself.
This is what makes \PoC{}'s privacy clause (Definition~\ref{def:poc}) more than a
promise: the contributor proves the gradient was computed over data consistent with
the committed policy, and shares only the resulting delta, so the corpus stays private
while the contribution stays verifiable.

\subsection{Weights are a portable, user-owned asset}

A model improvement---the $\Delta W$ of \PoT{}/\PoC{}---is serialized in the gym's
canonical backend-independent BF16 format and committed by its root
(Section~\ref{sec:four-proofs}). Because the serialization is canonical and
backend-independent, a delta trained on one accelerator is byte-portable to any other,
and its provenance (which base model, which data policy, what improvement margin) is
attested by the \PoC{} it was minted with. A user's weights are therefore an on-chain
asset they own and can move, sell, or compose---with royalties flowing to the
contributors whose proven improvements it incorporates---rather than an artifact
locked inside a single provider's stack. The proof framework is what gives the asset
its value: a delta is worth more when it carries a verifiable proof that it was
genuinely trained and genuinely improves the model, than when it is merely asserted to.

\subsection{Trust as a dial, security as a floor}

The synthesis is that the execution proofs decouple \emph{trust} from \emph{security}.
A user dials trust up the spectrum---from running everything themselves, to trusting a
proof system, to trusting a named operator---according to their convenience and
threat model, and the proof framework holds the security floor constant underneath:
in every mode, a result that claims to be the model's output on the user's input
either is, or is caught. This is the difference between the deployed settlement
engine, where ``trusting the operators'' meant accepting unverified guesses
(Section~\ref{sec:guess-to-mint}), and a proof-backed network, where trusting an
operator is a choice about \emph{speed and cost}, never about \emph{correctness}. The
user owns the data, owns the inference, and owns the weights; the protocol's job is
only to make all three auditable, at whatever point on the trust spectrum the user
selects.
